\documentclass[11pt,a4paper]{article}
\usepackage[utf8]{inputenc}
\usepackage[T1]{fontenc}
\usepackage[french]{babel}
\usepackage{amsmath,amssymb}
\usepackage{geometry}
\usepackage{enumitem}
\usepackage{booktabs}
\usepackage{array}
\usepackage{tabularx}
\usepackage{xcolor}
\usepackage{tcolorbox}
\usepackage{fancyhdr}
\usepackage{multicol}
\usepackage{tikz}
\usepackage{pgfplots}
\pgfplotsset{compat=1.18}
\usepackage{lastpage}

\tcbuselibrary{skins,breakable}

\geometry{margin=2cm, headheight=15pt}

% Couleurs
\definecolor{primary}{RGB}{31,78,121}
\definecolor{soft}{RGB}{230,239,247}
\definecolor{amberbg}{RGB}{255,244,224}
\definecolor{amberborder}{RGB}{240,192,112}
\definecolor{greenbg}{RGB}{233,246,234}
\definecolor{greenborder}{RGB}{182,224,184}
\definecolor{redbg}{RGB}{253,236,236}
\definecolor{redborder}{RGB}{242,179,171}
\definecolor{grayline}{RGB}{120,130,140}

% En-tête et pied de page
\pagestyle{fancy}
\fancyhf{}
\fancyhead[L]{\textcolor{primary}{\textbf{BTS BioALC — 2\textsuperscript{ème} année}}}
\fancyhead[R]{\textcolor{primary}{Fiche de travail — Statistique descriptive}}
\fancyfoot[L]{\textcolor{grayline}{\small Nom : \dotfill}}
\fancyfoot[C]{\textcolor{grayline}{\small \thepage/\pageref{LastPage}}}
\fancyfoot[R]{\textcolor{grayline}{\small Semaine du 05/10/2026}}
\renewcommand{\headrulewidth}{0.4pt}
\renewcommand{\footrulewidth}{0.4pt}

% Boîtes
\newtcolorbox{consigne}[1][]{
  colback=soft, colframe=primary, boxrule=0.8pt, arc=4pt,
  fonttitle=\bfseries, coltitle=white, colbacktitle=primary,
  left=6pt, right=6pt, top=4pt, bottom=4pt,
  breakable, #1
}

\newtcolorbox{contexte}{
  colback=amberbg, colframe=amberborder, boxrule=0.8pt, arc=4pt,
  title={\textbf{Contexte professionnel}}, fonttitle=\bfseries,
  coltitle=black, colbacktitle=amberborder,
  left=6pt, right=6pt, top=4pt, bottom=4pt,
  breakable
}

\newtcolorbox{rappel}{
  colback=greenbg, colframe=greenborder, boxrule=0.8pt, arc=4pt,
  title={\textbf{Rappel — Fiche méthode}}, fonttitle=\bfseries,
  coltitle=black, colbacktitle=greenborder,
  left=6pt, right=6pt, top=4pt, bottom=4pt,
  breakable
}

\newtcolorbox{questionbox}{
  colback=white, colframe=primary, boxrule=0.8pt, arc=4pt,
  fonttitle=\bfseries, coltitle=white, colbacktitle=primary,
  left=6pt, right=6pt, top=4pt, bottom=4pt,
  breakable
}

\newtcolorbox{autoeval}{
  colback=redbg, colframe=redborder, boxrule=0.8pt, arc=4pt,
  title={\textbf{Auto-évaluation}}, fonttitle=\bfseries,
  coltitle=black, colbacktitle=redborder,
  left=6pt, right=6pt, top=4pt, bottom=4pt,
  breakable
}

% Commande pour lignes d'écriture
\newcommand{\lignes}[1]{\vspace{0.15cm}\foreach \x in {1,...,#1}{\noindent\rule{\linewidth}{0.4pt}\par\vspace{0.35cm}}}

% Cases à cocher
\newcommand{\case}{\raisebox{0.1ex}{\fbox{\rule{0pt}{1.1ex}\rule{1.1ex}{0pt}}}}

\title{\textcolor{primary}{\textbf{Fiche de travail — Statistique descriptive}}\\
\large Semaine du 5 au 9 octobre 2026 — BTS Bioanalyses en Laboratoire de Contrôle}
\author{}
\date{}

\begin{document}

\maketitle

\begin{consigne}
\textbf{Mode d'emploi de cette fiche :} complétez directement sur ce document au fur et à mesure des deux séances. Les espaces libres sont prévus pour vos réponses, vos calculs et vos schémas. Cette fiche est un support de travail : elle n'est pas notée, mais elle vous servira pour réviser. En cas de difficulté, reportez-vous aux fiches méthodes rappelées dans chaque partie.
\end{consigne}

\vspace{0.3cm}

\begin{contexte}
Vous êtes technicien(ne) dans un laboratoire d'analyses médicales. Le responsable qualité vous demande d'analyser les résultats d'une série de mesures de \textbf{taux de glucose} (en g/L) réalisées sur \textbf{20 patients} dans le cadre d'un contrôle de routine. Votre mission : produire un tableau de synthèse, un graphique et les indicateurs statistiques permettant de valider la conformité de la série.
\end{contexte}

\vspace{0.4cm}

% ============================================================
\section*{Partie 1 — Séance du lundi (50 min)}
% ============================================================

\subsection*{1.1. Rebrassage — Questions orales}

\begin{questionbox}
Répondez brièvement aux trois questions suivantes :
\begin{enumerate}[leftmargin=*]
    \item Qu'appelle-t-on \textbf{population} en statistique ? Et \textbf{échantillon} ?
    \lignes{2}
    \item Quelle est la différence entre une variable \textbf{qualitative} et une variable \textbf{quantitative} ?
    \lignes{2}
    \item Qu'est-ce qu'un \textbf{effectif} ? Une \textbf{fréquence} ?
    \lignes{2}
\end{enumerate}
\end{questionbox}

\subsection*{1.2. Fiche méthode n\textsuperscript{o}1 — Représenter des données statistiques}

\begin{rappel}
\begin{tabularx}{\linewidth}{|p{3.5cm}|X|}
\hline
\textbf{Étape} & \textbf{À retenir} \\
\hline
1. Identifier la variable & Qualitative ou quantitative ? Discrète ou continue ? \\
\hline
2. Construire le tableau & Effectifs, fréquences, fréquences cumulées. \\
\hline
3. Choisir le graphique & Qualitative $\rightarrow$ diagramme en bâtons ou circulaire. \\
 & Quantitative discrète $\rightarrow$ diagramme en bâtons. \\
 & Quantitative continue $\rightarrow$ histogramme. \\
 & Cumul $\rightarrow$ courbe cumulative. \\
\hline
4. Vérifier & Titre, légende, unités, échelle correcte. \\
\hline
\end{tabularx}

\vspace{0.2cm}
\textbf{Exemple type :} taux de glucose (g/L) sur 20 patients $\rightarrow$ histogramme avec classes de 0,1 g/L.
\end{rappel}

\subsection*{1.3. Exercice 1 — Construire un tableau et un graphique}

\begin{questionbox}
Voici les taux de glucose (en g/L) mesurés sur 20 patients :

\begin{center}
\begin{tabular}{|c|c|c|c|c|c|c|c|c|c|}
\hline
0,82 & 0,91 & 0,75 & 1,02 & 0,88 & 0,95 & 0,79 & 1,10 & 0,85 & 0,93 \\
\hline
0,87 & 0,98 & 0,81 & 0,76 & 1,05 & 0,89 & 0,92 & 0,84 & 0,97 & 0,86 \\
\hline
\end{tabular}
\end{center}

\textbf{Question 1.} Regroupez ces données en \textbf{classes de 0,1 g/L} (de 0,70 à 1,10 g/L) et complétez le tableau ci-dessous.

\vspace{0.2cm}
\begin{center}
\begin{tabular}{|c|c|c|c|}
\hline
\textbf{Classe (g/L)} & \textbf{Effectif $n_i$} & \textbf{Fréquence $f_i$} & \textbf{Fréquence cumulée croissante} \\
\hline
[0,70 ; 0,80[ & & & \\
\hline
[0,80 ; 0,90[ & & & \\
\hline
[0,90 ; 1,00[ & & & \\
\hline
[1,00 ; 1,10] & & & \\
\hline
\textbf{Total} & & & \\
\hline
\end{tabular}
\end{center}

\vspace{0.3cm}
\textbf{Question 2.} Quel type de graphique est adapté à cette variable ? Justifiez.

\lignes{2}

\textbf{Question 3.} Construisez ce graphique dans le cadre ci-dessous (pensez à indiquer titre, axes et unités).

\vspace{0.2cm}
\begin{center}
\begin{tikzpicture}
\draw[gray!50, step=0.5cm] (0,0) grid (14,5);
\draw[thick, ->] (0,0) -- (14.3,0);
\draw[thick, ->] (0,0) -- (0,5.3);
\end{tikzpicture}
\end{center}
\end{questionbox}

\subsection*{1.4. Exercice 2 — Analyse critique de graphiques}

\begin{questionbox}
Deux graphiques vous sont présentés au tableau par l'enseignant.

\textbf{Question 1.} Quel est le graphique correct ? Pourquoi ?

\lignes{3}

\textbf{Question 2.} Listez au moins trois erreurs présentes dans le graphique incorrect.

\lignes{4}
\end{questionbox}

\subsection*{1.5. Synthèse de la séance 1 — Questions flash}

\begin{questionbox}
\begin{enumerate}[leftmargin=*]
    \item Quel graphique utiliser pour une variable qualitative ?
    \lignes{1}
    \item Que représente la fréquence cumulée croissante ?
    \lignes{1}
    \item Pourquoi faut-il toujours indiquer les unités sur un graphique ?
    \lignes{1}
\end{enumerate}
\end{questionbox}

\newpage

% ============================================================
\section*{Partie 2 — Séance du jeudi (50 min)}
% ============================================================

\subsection*{2.1. Rebrassage — Rappel de la séance 1}

\begin{questionbox}
\begin{enumerate}[leftmargin=*]
    \item Qu'est-ce qu'un tableau d'effectifs ?
    \lignes{1}
    \item Quelle est la différence entre un diagramme en bâtons et un histogramme ?
    \lignes{2}
\end{enumerate}
\end{questionbox}

\subsection*{2.2. Fiche méthode n\textsuperscript{o}2 — Moyenne, médiane, mode}

\begin{rappel}
\begin{tabularx}{\linewidth}{|p{2.4cm}|X|p{3.8cm}|}
\hline
\textbf{Indicateur} & \textbf{Définition / Formule} & \textbf{Interprétation} \\
\hline
\textbf{Moyenne} & $\bar{x} = \dfrac{1}{n}\displaystyle\sum_{i=1}^{n} x_i$ & Sensible aux valeurs extrêmes. \\
\hline
\textbf{Médiane} & Valeur qui partage la série en deux groupes égaux (trier puis prendre la valeur centrale, ou la moyenne des deux valeurs centrales). & Robuste aux valeurs extrêmes. \\
\hline
\textbf{Mode} & Valeur la plus fréquente (compter les effectifs). & Utile pour les variables qualitatives. \\
\hline
\end{tabularx}

\vspace{0.2cm}
\textbf{Exemple type :} série $2, 3, 3, 4, 5, 5, 5, 6, 7, 8$.

Moyenne : $\bar{x} = \dfrac{2+3+3+4+5+5+5+6+7+8}{10} = \dfrac{48}{10} = 4{,}8$.

Médiane : effectif pair ($n=10$), donc moyenne des 5\textsuperscript{e} et 6\textsuperscript{e} valeurs triées : $\dfrac{5+5}{2} = 5$.

Mode : la valeur la plus fréquente est 5 (effectif 3).
\end{rappel}

\subsection*{2.3. Exercice 1 — Calculer moyenne, médiane et mode}

\begin{questionbox}
Reprenez les données de glucose (20 patients) et répondez aux questions suivantes.

\textbf{Question 1.} Complétez le tableau de calcul de la moyenne (utilisation des centres de classes).

\vspace{0.2cm}
\begin{center}
\begin{tabular}{|c|c|c|c|}
\hline
\textbf{Classe (g/L)} & \textbf{Centre $c_i$} & \textbf{Effectif $n_i$} & \textbf{$n_i \times c_i$} \\
\hline
[0,70 ; 0,80[ & 0,75 & & \\
\hline
[0,80 ; 0,90[ & 0,85 & & \\
\hline
[0,90 ; 1,00[ & 0,95 & & \\
\hline
[1,00 ; 1,10] & 1,05 & & \\
\hline
\textbf{Total} & & & \\
\hline
\end{tabular}
\end{center}

\vspace{0.2cm}
Calculez la moyenne : $\bar{x} = \dfrac{\dots\dots\dots}{\dots\dots\dots} = \dots\dots\dots$ g/L

\vspace{0.3cm}
\textbf{Question 2.} Triez les 20 valeurs brutes par ordre croissant, puis déterminez la médiane.

\vspace{0.2cm}
Valeurs triées : \dotfill

\vspace{0.2cm}
Valeurs triées (suite) : \dotfill

\vspace{0.2cm}
Médiane : \dotfill

\vspace{0.3cm}
\textbf{Question 3.} Quel est le mode de cette série ? Justifiez.

\lignes{2}

\textbf{Question 4.} Comparez la moyenne et la médiane. Que peut-on dire de la distribution ?

\lignes{3}
\end{questionbox}

\subsection*{2.4. Exercice 2 — Utilisation du tableur}

\begin{questionbox}
\textbf{Étape 1.} Ouvrez un tableur et saisissez les 20 valeurs brutes dans la colonne A (de A1 à A20).

\textbf{Étape 2.} Dans les cellules indiquées, saisissez les formules suivantes :

\vspace{0.2cm}
\begin{center}
\begin{tabular}{|c|c|l|}
\hline
\textbf{Cellule} & \textbf{Formule à saisir} & \textbf{Résultat obtenu} \\
\hline
C1 & \texttt{=MOYENNE(A1:A20)} & \dotfill \\
\hline
C2 & \texttt{=MEDIANE(A1:A20)} & \dotfill \\
\hline
C3 & \texttt{=MODE(A1:A20)} & \dotfill \\
\hline
\end{tabular}
\end{center}

\vspace{0.3cm}
\textbf{Étape 3.} Comparez les résultats obtenus avec ceux de l'exercice 1. Sont-ils identiques ?

\lignes{2}

\textbf{Étape 4.} Construisez un diagramme en bâtons des effectifs. Ajoutez une ligne horizontale représentant la moyenne. Décrivez ce que vous observez.

\lignes{3}
\end{questionbox}

\subsection*{2.5. Évaluation formative — QCM}

\begin{questionbox}
\textbf{QCM de 5 questions (sans note — correction immédiate en classe).} Entourez la bonne réponse.

\vspace{0.3cm}
\textbf{Question 1.} La moyenne est un indicateur :

\hspace{1cm} A. de position \hspace{2cm} B. de dispersion \hspace{2cm} C. de forme

\vspace{0.3cm}
\textbf{Question 2.} Pour la série $4, 5, 5, 6, 7, 7, 7, 8, 9, 10$, la médiane vaut :

\hspace{1cm} A. 5 \hspace{2.5cm} B. 6{,}5 \hspace{2.5cm} C. 7

\vspace{0.3cm}
\textbf{Question 3.} Le mode de la série $2, 3, 3, 4, 4, 4, 5$ est :

\hspace{1cm} A. 2 \hspace{2.5cm} B. 4 \hspace{2.5cm} C. 5

\vspace{0.3cm}
\textbf{Question 4.} Un écart important entre moyenne et médiane indique :

\hspace{1cm} A. une distribution symétrique

\hspace{1cm} B. la présence de valeurs extrêmes

\hspace{1cm} C. une erreur de calcul

\vspace{0.3cm}
\textbf{Question 5.} Dans un contexte de contrôle qualité, la médiane est souvent préférée à la moyenne car :

\hspace{1cm} A. elle est plus facile à calculer

\hspace{1cm} B. elle est moins sensible aux valeurs aberrantes

\hspace{1cm} C. elle est toujours plus grande
\end{questionbox}

\subsection*{2.6. Synthèse de la séance 2}

\begin{questionbox}
\textbf{À retenir :} complétez les phrases suivantes.

\begin{itemize}[leftmargin=*]
    \item La moyenne est sensible aux \dotfill
    \item La médiane est \dotfill aux valeurs extrêmes.
    \item Le mode est la valeur \dotfill de la série.
    \item Un écart important entre moyenne et médiane signale la présence de \dotfill
\end{itemize}
\end{questionbox}

\newpage

% ============================================================
\section*{Bilan de la semaine}
% ============================================================

\begin{autoeval}
Cochez la case correspondant à votre niveau pour chaque compétence :

\vspace{0.3cm}
\begin{center}
\begin{tabularx}{\linewidth}{|X|c|c|c|}
\hline
\textbf{Compétence} & \textbf{Acquis} & \textbf{En cours} & \textbf{À revoir} \\
\hline
Je sais construire un tableau d'effectifs et de fréquences. & \case & \case & \case \\
\hline
Je sais choisir un graphique adapté. & \case & \case & \case \\
\hline
Je sais construire un histogramme. & \case & \case & \case \\
\hline
Je sais calculer une moyenne. & \case & \case & \case \\
\hline
Je sais déterminer une médiane. & \case & \case & \case \\
\hline
Je sais identifier le mode. & \case & \case & \case \\
\hline
Je sais utiliser les fonctions du tableur. & \case & \case & \case \\
\hline
Je sais interpréter les indicateurs dans un contexte de laboratoire. & \case & \case & \case \\
\hline
\end{tabularx}
\end{center}

\vspace{0.3cm}
\textbf{Une question que je me pose encore :}

\lignes{2}
\end{autoeval}

\vspace{0.4cm}

\begin{contexte}
\textbf{Pour aller plus loin (facultatif, non noté) :} recherchez dans une norme ISO (par exemple ISO 5725) comment la moyenne et l'écart-type sont utilisés pour évaluer la \textbf{fidélité} d'une méthode de mesure. Rédigez un court paragraphe (5 lignes) expliquant l'intérêt de ces indicateurs en contrôle qualité.

\lignes{5}
\end{contexte}

\vspace{0.5cm}

\begin{center}
\textit{Fiche de travail — BTS BioALC — 2\textsuperscript{ème} année — Semaine du 05/10/2026}\\
\textit{Document à conserver pour les révisions.}
\end{center}

\end{document}